\documentclass[twoside,12pt]{book}

\usepackage[utf8]{inputenc}
\usepackage[T1]{fontenc}
\usepackage{mathpazo}
\usepackage{amsmath,amssymb,amsthm}
\usepackage{graphicx}
\usepackage{listings}
\usepackage{xcolor}
\usepackage{enumitem}
\usepackage{tikz}
\usepackage{algorithm}
\usepackage{algpseudocode}
\usepackage{booktabs}
\usepackage{geometry}
\usepackage{fancyhdr}
\usepackage{titlesec}
\usepackage[colorlinks=true, linkcolor=blue, citecolor=blue, urlcolor=blue, plainpages=false, pdfpagelabels]{hyperref}

\geometry{
  papersize={6in,9in},
  margin=0.75in,
  top=0.75in,bottom=0.75in,
  inner=0.85in,outer=0.65in,
  headheight=14pt
}

\pagestyle{fancy}
\fancyhf{}
\fancyhead[LE]{\small\textit{2. Fundamentals}}
\fancyhead[RO]{\small\textit{Randomized Algorithms}}
\fancyfoot[C]{\thepage}
\renewcommand{\headrulewidth}{0.4pt}

\definecolor{codegray}{rgb}{0.45,0.45,0.45}
\lstdefinestyle{cppstyle}{
  basicstyle=\ttfamily\scriptsize,
  keywordstyle=\bfseries,
  commentstyle=\itshape\color{codegray},
  numbers=left,
  numberstyle=\tiny\color{codegray},
  numbersep=8pt,
  frame=single,
  framerule=0.4pt,
  rulecolor=\color{codegray},
  backgroundcolor=\color{gray!5},
  breaklines=true,
  tabsize=4,
  showstringspaces=false,
  language=C++,
  morekeywords={nullptr,size_t,auto,const,constexpr,
    unordered_map,vector,pair,mt19937,
    uniform_int_distribution,INT_MAX}
}
\lstset{style=cppstyle}

\theoremstyle{plain}
\newtheorem{theorem}{Theorem}[chapter]
\newtheorem{lemma}[theorem]{Lemma}
\newtheorem{corollary}[theorem]{Corollary}
\newtheorem{proposition}[theorem]{Proposition}
\theoremstyle{definition}
\newtheorem{definition}[theorem]{Definition}
\newtheorem{example}[theorem]{Example}
\newtheorem{remark}[theorem]{Remark}
\theoremstyle{remark}
\newtheorem{exercise}[theorem]{Exercise}
\newtheorem{problem}[theorem]{Problem}
\newenvironment{cpplisting}{\begin{center}\small}{\end{center}}

\newcommand{\Exp}{\operatorname{\mathbb{E}}}
\newcommand{\Var}{\operatorname{Var}}
\newcommand{\Cov}{\operatorname{Cov}}

\begin{document}

\chapter*{2 Fundamentals}
\addcontentsline{toc}{chapter}{2 Fundamentals}
\markboth{2 Fundamentals}{Randomized Algorithms}
\label{ch:fundamentals}

In this chapter we develop fundamental mathematical tools and computational
models that underpin the randomized algorithms studied in the remainder of
the book. We begin with the construction of binary planar partitions,
illustrating linearity of expectation with a geometric application. We then
study a probabilistic recurrence that arises in the analysis of selection
and other divide-and-conquer algorithms. Finally, we review the computation
model and complexity classes that provide the formal framework for our
discussions throughout the book.

\section{Binary Planar Partitions}
\label{sec:bpp}
%==========================================================================

We now illustrate linearity of expectation with another application. Given a
set $S=\{s_1,s_2,\ldots,s_n\}$ of non-intersecting line segments in the
plane, we wish to construct a \emph{binary planar partition}---a binary
tree in which each node is associated with a region of the plane, and each
region contains at most one line segment. Such partitions have applications
to hidden-line elimination in computer graphics.

%% ----- Binary Planar Partitions -----
\subsection{Binary Planar Partitions}

Our next illustration is the construction of a binary planar partition of a
set of $n$ disjoint line segments in the plane, a problem with applications
to computer graphics.  A binary planar partition consists of a binary tree
together with some additional information, as described below.  Every internal
node of the tree has two children.  Associated with each node $v$ of the tree
is a region $r(v)$ of the plane.  Associated with each internal node $v$ of
the tree is a line $\ell(v)$ that intersects $r(v)$.  The region
corresponding to the root is the entire plane.  The region $r(v)$ is
partitioned by $\ell(v)$ into two regions $r_1(v)$ and $r_2(v)$, which are
the regions associated with the two children of~$v$.  Thus, any region $r$ of
the partition is bounded by the partition lines on the path from the root to
the node corresponding to $r$ in the tree.

Given a set $S=\{s_1,s_2,\ldots,s_n\}$ of non-intersecting line segments
in the plane, we wish to find a binary planar partition such that every
region in the partition contains at most one line segment (or a portion of
one line segment).  Notice that the definition allows us to divide an input
line segment $s_i$ into several segments $s_{i1},s_{i2},\ldots$, each of
which lies in a different region.

\begin{exercise}
Show that there exists a set of line segments for which no binary planar
partition can avoid breaking up some of the segments into pieces, if each
segment is to lie in a different region of the partition.
\end{exercise}

Binary planar partitions have two applications in computer graphics.  Here,
we describe one of them, the problem of hidden line elimination.  The second
application has to do with the constructive solid geometry (or CSG)
representation of a polyhedral object.

In rendering a scene on a graphics terminal, we are often faced with a
situation in which the scene remains fixed, but it is to be viewed from
several directions.  The hidden line elimination problem is the following:
having adopted a viewpoint and a direction of viewing, we want to draw only
the portion of the scene that is visible, eliminating those objects that are
``obscured by other objects'' in front of them relative to the viewpoint.

One approach uses a binary partition tree.  Once the scene has been decomposed
into line segments, we construct a binary planar partition tree for it.  Now,
given the direction of viewing, we use an idea known as the \emph{painter's
algorithm} to render the scene: first draw the objects that are furthest
``behind,'' and then progressively draw the objects that are in front.

The time it takes to render the scene depends on the size of the binary
planar partition tree.  We therefore wish to construct a binary planar
partition that is as small as possible.  Because the construction of the
partition can break some of the input segments $s_i$ into smaller pieces, the
size of the partition need not be~$n$; in fact, it is not clear that a
partition of size $O(n)$ always exists.

For a line segment $s$, let $\ell(s)$ denote the line obtained by extending
(if necessary) $s$ on both sides to infinity.  For the set
$S=\{s_1,s_2,\ldots,s_n\}$ of line segments, a simple and natural class of
partitions is the set of \emph{autopartitions}, which are formed by only
using lines from the set $\{\ell(s_1),\ell(s_2),\ldots,\ell(s_n)\}$ in
constructing the partition.

\begin{algorithm}[htbp]
\caption{RandAuto}
\label{alg:randauto}
\begin{algorithmic}[1]
\State \textbf{Input:} A set $S=\{s_1,\ldots,s_n\}$ of non-intersecting
  line segments.
\State \textbf{Output:} A binary autopartition $P_\pi$ of $S$.
\State Pick a permutation $\pi$ of $\{1,2,\ldots,n\}$ uniformly at random.
\While{a region contains more than one segment}
  \State Cut it with $\ell(s_i)$ where $i$ is first in $\pi$ such that
    $s_i$ cuts that region.
\EndWhile
\end{algorithmic}
\end{algorithm}

\begin{theorem}
\label{thm:partition}
The expected size of the autopartition produced by RandAuto is $O(n\log n)$.
\end{theorem}

\begin{proof}
For line segments $u$ and $v$, define $\operatorname{index}(u,v)$ to be~$i$
if $\ell(u)$ intersects $i{-}1$ other segments before hitting $v$, and
$\operatorname{index}(u,v)=\infty$ if $\ell(u)$ does not hit $v$.  Since a
segment $u$ can be extended in two directions, it is possible that
$\operatorname{index}(u,v)=\operatorname{index}(u,w)$ for two different
lines $v$ and $w$.

Let us denote by $u\dashv v$ the event that $\ell(u)$ cuts $v$ in the
constructed partition.  Let $\operatorname{index}(u,v)=i$, and let
$u_1,u_2,\ldots,u_{i-1}$ be the segments that $\ell(u)$ intersects before
hitting~$v$.  The event $u\dashv v$ happens only if $u$ occurs before any
of $\{u_1,u_2,\ldots,u_{i-1},v\}$ in the randomly chosen permutation
$\pi$.  The probability that this happens is $1/(i+1)$.

Let $C_{u,v}$ be an indicator variable that is 1 if $u\dashv v$ and 0
otherwise; clearly,
$\mathbb{E}[C_{u,v}]=\Pr[u\dashv v]\le 1/(\operatorname{index}(u,v)+1)$.
The size of $P_\pi$ equals $n$ plus the number of intersections due to cuts.
Thus, its expectation is
$n+\mathbb{E}[\sum_u\sum_v C_{u,v}]$ and by linearity of expectation this
equals
\begin{equation*}
  n + \sum_u \sum_v \Pr[u\dashv v]
  \le n + \sum_u \sum_v
    \frac{1}{\operatorname{index}(u,v)+1}.
  \tag{2.1}
\end{equation*}

For any line segment $u$ and any finite positive integer $i$, there are at
most two vertices $v$ and $w$ such that $\operatorname{index}(u,v)$ and
$\operatorname{index}(u,w)$ equal~$i$.  This is because the extension of the
segment $u$ along either of the two possible directions will meet any other
line segment at most once.  Thus, in each of the two directions, there is a
total ordering on the points of intersection with other segments and the
index values increase monotonically.  This implies that
\begin{equation*}
  \sum_{v\ne u} \frac{1}{\operatorname{index}(u,v)+1}
  \le \sum_{i=1}^{n-1} \frac{2}{i+1}
  = O(\log n).
  \tag{2.2}
\end{equation*}

Combining this with~(2.1) implies that the expected size of $P_\pi$ is
bounded above by $n+O(n\log n)$, which is $O(n\log n)$.
\end{proof}

Note that in computing the expected number of intersections, we only made
use of linearity of expectation.  We do not require any independence between
the events $u\dashv v$ and $u\dashv w$, for segments $u$, $v$, and~$w$.
Indeed, these events need not be independent in general.

One way of interpreting Theorem~\ref{thm:partition} is as follows: since the
expected size of the binary planar partition constructed by the algorithm is
$O(n\log n)$ on any input, there must exist a binary autopartition of size
$O(n\log n)$ for every input.  This follows from the simple fact that any
random variable assumes at least one value that is no greater than its
expectation.  Thus we have used a \emph{probabilistic argument} to assert
that a combinatorial object---in this case a binary autopartition of size
$O(n\log n)$---exists with absolute certainty rather than with some
probability.  This is an example of the \emph{probabilistic method} in
combinatorics.  We will study the probabilistic method in greater detail in
Chapter~6.

%% ----- C++ : BPP -----
\subsubsection*{C++ Implementation}

\begin{lstlisting}[
  caption={RandAuto for binary planar partitions.},
  label={lst:bpp}]
struct Seg { double x1,y1,x2,y2; int id; };

bool intersect(const Seg& a, const Seg& b) {
  auto o=[](double px,double py,
            double qx,double qy,
            double rx,double ry){
    double v=(qy-py)*(rx-qx)-(qx-px)*(ry-qy);
    return std::abs(v)<1e-10?0:(v>0?1:2);
  };
  return o(a.x1,a.y1,a.x2,a.y2,b.x1,b.y1)
        !=o(a.x1,a.y1,a.x2,a.y2,b.x2,b.y2)
      && o(b.x1,b.y1,b.x2,b.y2,a.x1,a.y1)
        !=o(b.x1,b.y1,b.x2,b.y2,a.x2,a.y2);
}

// Build autopartition recursively
// segs: remaining segments in this region
// Returns number of leaf nodes (partition size)
int rand_auto(std::vector<Seg*>& segs) {
  if (segs.empty()) return 0;
  if (segs.size()==1) return 1;       // leaf
  std::vector<int> p(segs.size());
  std::iota(p.begin(),p.end(),0);
  std::shuffle(p.begin(),p.end(),
    std::mt19937(std::random_device()()));
  Seg* cut = segs[p[0]];             // random pivot
  std::vector<Seg*> L,R;
  for (size_t i=1;i<segs.size();i++) {
    if (intersect(*cut,*segs[i]))
      { L.push_back(segs[i]); R.push_back(segs[i]); }
    else {
      double dx=cut->x2-cut->x1;
      double dy=cut->y2-cut->y1;
      double cross=dx*(segs[i]->y1-cut->y1)
                  -dy*(segs[i]->x1-cut->x1);
      (cross>=0?L:R).push_back(segs[i]);
    }
  }
  return rand_auto(L)+rand_auto(R);
}
\end{lstlisting}

%==========================================================================
\section{A Probabilistic Recurrence}
\label{sec:recurrence}
%==========================================================================

Frequently, we express a random variable of interest as a recurrence in
terms of other random variables.  In this section, we study one such
situation using the Find algorithm analyzed in detail in Problem~1.9.  The
material in this section, although useful, is not an essential prerequisite
for subsequent topics and may be omitted in the first reading.

The Find algorithm for selecting the $k$th smallest of a set $S$ of $n$
elements works as follows.  We pick a random element $y$ and partition
$S\setminus\{y\}$ into two sets $S_1$ and $S_2$ (elements smaller and larger
than~$y$ respectively) as in RandQS.  Suppose $|S_1|=k-1$; then $y$ is the
desired element and we are done.  Otherwise, if $|S_1|>k$, we recursively
find the $k$th smallest element of $S_1$; else we recursively find the
$(k-|S_1|-1)$th smallest element in~$S_2$.

The expected number of comparisons made by the Find algorithm is the subject
of Problem~1.9.  Suppose instead that we were to ask the following question:
what is the expected number of times we make the recursive call in the
algorithm?  Equivalently, what is the expected number of times we pick a
random element in the algorithm?  While this question may not be especially
important for the Find algorithm, it is the kind of question that arises in
the analysis of a number of parallel and geometric algorithms.  Intuitively,
we expect that the size of the residual problem in the Find algorithm is
divided by a constant factor at each recursive level, so that we expect that
the number of recursive invocations is $O(\log n)$.  Below, we show that
this intuition can be formalized in a general setting.

Let $g(x)$ be a monotone non-decreasing function from the positive reals to
the positive reals.  Consider a particle whose position changes at discrete
time steps and is always at a positive integer.  If the particle is currently
at position $m>1$, it proceeds at the next step to the position $m-X$, where
$X$ is a random variable ranging over the integers $1,\ldots,m-1$.  All we
know about $X$ is that $\mathbb{E}[X]\ge g(m)$, and that $X$ is chosen
independently of the past.  It is clear that the particle will always reach
position~1 and the process terminates in that state.  The interesting
question is, assuming that the particle starts at position~$n$, what is the
expected number of steps before it reaches position~1?

\begin{theorem}
\label{thm:recurrence}
Let $T$ be the random variable denoting the number of steps in which the
particle reaches the position~1.  Then,
$\mathbb{E}[T]\le \displaystyle\int_1^n \frac{dx}{g(x)}$.
\end{theorem}

\begin{proof}
The proof is by induction on $n$; let us suppose the theorem holds for values
of $m$ smaller than $n$.  Let $f(m)=\int_1^m dx/g(x)$ for $m>1$.  We wish
to show that $\mathbb{E}[T]\le f(n)$.

Consider the first step, during which the particle proceeds from position $n$
to position $n-X$, where $X$ is chosen from a distribution for which
$\mathbb{E}[X]\ge g(n)$.  We have
\begin{align}
  \mathbb{E}[T]
  &\le 1 + \mathbb{E}[f(n-X)]
    \tag{1.9}\\
  &= 1 + \mathbb{E}\!\left[\int_{n-X}^{n}\frac{dy}{g(y)}\right]
    \tag{1.10}\\
  &\le 1 + \mathbb{E}\!\left[\frac{X}{g(n)}\right]
    \tag{1.11}\\
  &= 1 + \frac{\mathbb{E}[X]}{g(n)}
    \tag{1.12}\\
  &\le 1 + f(n) - 1 = f(n).
    \tag{1.13}
\end{align}
The inequality~(1.10) follows from the inductive hypothesis applied to
position $n-X$.  The inequality~(1.11) follows from the assumption that $g$
is non-decreasing, so $1/g(y)\le 1/g(n)$ for $y\le n$.  The
inequality~(1.13) follows from the lower bound on $\mathbb{E}[X]$ and the
observation that $f(n)-f(n-X)\ge X/g(n)$.
\end{proof}

\begin{exercise}
If $X$ were to range over all integers having value at most $m-1$ (possibly
including negative integers), how would the statement and proof of
Theorem~\ref{thm:recurrence} change?
\end{exercise}

For the Find algorithm, we can show (following the analysis of Problem~1.9)
that $g(m)\ge m/4$.  We may then apply the above theorem to bound the
expected number of recursive calls to Find by $4\ln n$.

\begin{exercise}
What prevents us from using Theorem~\ref{thm:recurrence} to bound the
expected number of levels of recursion in the RandQS algorithm?
\end{exercise}

%% ----- C++ : Find -----
\subsection{C++ Implementation}

\begin{lstlisting}[
  caption={Randomized Find---selecting the $k$th smallest element.},
  label={lst:find}]
int rand_find(std::vector<int>& a, int k) {
  if (a.size()==1) return a[0];
  int pi = RandomEngine::instance()
               .uniform_int(0,a.size()-1);
  int pivot = a[pi];
  std::vector<int> L,E,G;
  for (int x : a) {
    if      (x<pivot) L.push_back(x);
    else if (x==pivot) E.push_back(x);
    else              G.push_back(x);
  }
  if (k<=(int)L.size())
    return rand_find(L,k);
  else if (k<=(int)(L.size()+E.size()))
    return pivot;
  else
    return rand_find(G,k-L.size()-E.size());
}
\end{lstlisting}

\begin{lstlisting}[
  caption={Empirical measurement of recursion depth.},
  label={lst:find_analysis}]
struct Stats { int calls=0, depth=0; };

int rand_find_analyzed(std::vector<int>& a,int k,
                       int d, Stats& s) {
  s.calls++;
  s.depth = std::max(s.depth,d);
  if (a.size()==1) return a[0];
  int pi = RandomEngine::instance()
               .uniform_int(0,a.size()-1);
  int pivot = a[pi];
  std::vector<int> L,E,G;
  for (int x : a) {
    if      (x<pivot) L.push_back(x);
    else if (x==pivot) E.push_back(x);
    else              G.push_back(x);
  }
  if (k<=(int)L.size())
    return rand_find_analyzed(L,k,d+1,s);
  else if (k<=(int)(L.size()+E.size()))
    return pivot;
  else
    return rand_find_analyzed(G,
      k-L.size()-E.size(),d+1,s);
}

// Typical output for n=1000:
//   Avg recursive calls: ~12
//   Avg max depth:       ~11
//   Theoretical O(log n): ~10
//   Theorem 1.3 bound 4 ln n: ~27
\end{lstlisting}

%==========================================================================
\section{Computation Model and Complexity Classes}
\label{sec:complexity}
%==========================================================================

\subsection{RAMs and Turing Machines}

Following common practice, throughout this book we use the Turing machine
model to discuss complexity-theory issues.  As is common, however, we switch
to the RAM (random access machine) as the model of computation when
describing and analyzing algorithms (except in the study of parallel and
distributed algorithms in Chapter~13, where we define a version of the RAM
model for machines working in parallel).

\begin{definition}
A \emph{deterministic Turing machine} is a quadruple $M=(S,\Sigma,\delta,s)$.
Here $S$ is a finite set of states, of which $s\in S$ is the machine's
initial state.  The machine uses a finite set of symbols, denoted~$\Sigma$;
this set includes special symbols $\mathrm{BLANK}$ and $\mathrm{FIRST}$.  The
function $\delta$ is the transition function, mapping
$S\times\Sigma$ to
$(S\cup\{\mathrm{HALT},\mathrm{YES},\mathrm{NO}\})\times\Sigma\times\{\leftarrow,
\rightarrow,\mathrm{STAY}\}$.
\end{definition}

A \emph{probabilistic Turing machine} is a Turing machine augmented with the
ability to generate an unbiased coin flip in one step.  It corresponds to a
randomized algorithm.

In the RAM model, we have a machine that can perform the following types of
operations involving registers and main memory: input--output operations,
memory--register transfers, indirect addressing, branching, and arithmetic
operations.  Each register or memory location may hold an integer that can be
accessed as a unit, but an algorithm has no access to the representation of
the number.

\begin{proposition}
Any Turing machine computation of length polynomial in the size of the input
can be simulated by a RAM computation of length polynomial in the size of the
input.  Any RAM computation of length polynomial in the size of the input can
be simulated by a Turing machine computation of length polynomial in the size
of the input.
\end{proposition}

\subsection{Complexity Classes}

We now define some basic complexity classes focusing on those involving
randomized algorithms.

\begin{definition}
The class $\mathbf{P}$ consists of all languages $L$ that have a
polynomial-time algorithm $A$ such that for any input $x\in\Sigma^*$,
$x\in L \Rightarrow A(x)$ accepts, and
$x\notin L \Rightarrow A(x)$ rejects.
\end{definition}

\begin{definition}
The class $\mathbf{NP}$ consists of all languages $L$ that have a
polynomial-time algorithm $A$ such that for any input $x\in\Sigma^*$,
$x\in L \Rightarrow \exists\, y\in\Sigma^*: A(x,y)$ accepts, where
$|y|$ is bounded by a polynomial in $|x|$; and
$x\notin L \Rightarrow \forall\, y\in\Sigma^*: A(x,y)$ rejects.
\end{definition}

\begin{definition}
The class $\mathbf{RP}$ (Randomized Polynomial time) consists of all
languages $L$ that have a randomized algorithm $A$ running in worst-case
polynomial time such that for any input $x\in\Sigma^*$,
$x\in L \Rightarrow \Pr[A(x)\text{ accepts}]\ge 1/2$; and
$x\notin L \Rightarrow \Pr[A(x)\text{ accepts}]=0$.
\end{definition}

\begin{definition}
The class $\mathbf{ZPP}$ (Zero-error Probabilistic Polynomial time) is the
class of languages that have Las Vegas algorithms running in expected
polynomial time.
\end{definition}

\begin{definition}
The class $\mathbf{BPP}$ (Bounded-error Probabilistic Polynomial time)
consists of all languages $L$ that have a randomized algorithm $A$ running
in worst-case polynomial time such that for any input $x\in\Sigma^*$,
$x\in L \Rightarrow \Pr[A(x)\text{ accepts}]\ge 3/4$; and
$x\notin L \Rightarrow \Pr[A(x)\text{ accepts}]\le 1/4$.
\end{definition}

\begin{exercise}
Show that $\mathrm{ZPP}=\mathrm{RP}\cap\mathrm{co\text{-}RP}$.
\end{exercise}



%==========================================================================
\addcontentsline{toc}{section}{Notes}
\section*{Notes}
%==========================================================================

The ideas underlying randomized algorithms can be traced back to Monte Carlo
methods used in numerical analysis, statistical physics, and simulation.  In
the context of computability theory, the notion of a probabilistic Turing
machine was proposed by de~Leeuw, Moore, Shannon, and Shapiro~[122] and
further explored in the pioneering work of Rabin~[340] and
Gill~[166].  Berlekamp~[57], Rabin~[341], and Solovay and
Strassen~[382] gave early examples of concrete randomized algorithms.

Our RandQS algorithm is based on Hoare's algorithm~[201].  The min-cut
algorithm of Section~1.2.2, together many variations and
extensions, is due to Karger~[231].  Monte Carlo methods have been popular in
the sciences for over a hundred years.  The classic experiment on
approximating the value of $\pi$ by dropping needles on a sheet of paper with
parallel lines is described in an eighteenth-century paper by
Buffon~[86].  The origin of the modern theory of Monte Carlo methods in the
physical sciences is widely attributed to Ulam, von~Neumann, and
Fermi~[116].  The term \emph{Las Vegas algorithm} was introduced by
Babai~[37].

The material of Section~\ref{sec:bpp} is drawn from Paterson and
Yao~[329].  The Find algorithm described in Section~\ref{sec:recurrence} is
due to Hoare~[200].  Theorem~\ref{thm:recurrence} is given in a paper by
Karp, Upfal, and Wigderson~[250].  Karp~[244] gives a number of additional
results on probabilistic recurrence relations.

The reader is referred to introductory texts on algorithms and complexity
such as those by Aho, Hopcroft, and Ullman~[5,~6] and
Papadimitriou~[326] for more details on the Turing machine model and the
RAM model.  The books by Bovet and Crescenzi~[81] and by
Papadimitriou~[326] contain a more detailed treatment of the complexity
classes described in this chapter.

%==========================================================================
\addcontentsline{toc}{section}{Problems}
\section*{Problems}
%==========================================================================

\begin{exercise}[1.1]
(Due to J.~von~Neumann~[409].)
\begin{enumerate}[label=(\alph*)]
\item Suppose you are given a coin for which the probability of HEADS, say
  $p$, is unknown.  How can you use this coin to generate unbiased (i.e.,
  $\Pr[\text{HEADS}]=\Pr[\text{TAILS}]=1/2$) coin-flips?  Give a scheme for
  which the expected number of flips of the biased coin for extracting one
  unbiased coin-flip is no more than $1/[p(1-p)]$.  (Hint: Consider two
  consecutive flips of the biased coin.)
\item Devise an extension of the scheme that extracts the largest possible
  number of independent, unbiased coin-flips from a given number of flips of
  the biased coin.
\end{enumerate}
\end{exercise}

\begin{exercise}[1.2]
(Due to D.~E.~Knuth and A.~C.-C.~Yao~[264].)
\begin{enumerate}[label=(\alph*)]
\item Suppose you are provided with a source of unbiased random bits.  Explain
  how you will use this to generate uniform samples from the set
  $S=\{0,\ldots,n-1\}$.  Determine the expected number of random bits
  required by your sampling algorithm.
\item What is the worst-case number of random bits required by your sampling
  algorithm?  Consider the case when $n$ is a power of~2, as well as the case
  when it is not.
\item Solve (a) and (b) when, instead of unbiased random bits, you are
  required to use as the source of randomness uniform random samples from the
  set $\{0,\ldots,p-1\}$; consider the case when $n$ is a power of $p$, as
  well as the case when it is not.
\end{enumerate}
\end{exercise}

\begin{exercise}[1.3]
(Due to D.~E.~Knuth and A.~C.-C.~Yao~[264].) Suppose you are provided with
a source of unbiased random bits.  Provide efficient (in terms of expected
running time and expected number of random bits used) schemes for generating
samples from the distribution over the set $\{2,3,\ldots,12\}$ induced by
rolling two unbiased dice and taking the sum of their outcomes.
\end{exercise}

\begin{exercise}[1.4]
\begin{enumerate}[label=(\alph*)]
\item Suppose you are required to generate a random permutation of size $n$.
  Assuming that you have access to a source of independent and unbiased
  random bits, suggest a method for generating random permutations of size
  $n$.  Efficiency is measured in terms of both time and number of random
  bits.  What lower bounds can you prove for this task?
\item Consider the following method for generating a random permutation of
  size $n$.  Pick $n$ random values $X_1,\ldots,X_n$ independently from the
  uniform distribution over the interval $[0,1]$.  Now, the permutation that
  orders the random variables in ascending order is claimed to be a random
  permutation, and it can be determined by sorting the random values.  Is the
  claim correct?  How efficient is this scheme?
\item Consider the following ``lazy'' implementation of the scheme suggested
  in~(b).  The binary representation of the fraction $X_i$ is a sequence of
  unbiased and independent random bits.  At any given stage of the sorting
  algorithm, we would have chosen only as many bits of each $X_i$ as
  necessary to resolve all the comparisons performed up to that point.  When
  comparing $X_i$ to $X_j$, if the current prefixes of their binary
  expansions do not determine the outcome of the comparison, then we extend
  their prefixes by choosing further random bits until this happens.  Compute
  tight bounds on the expected number of random bits used by this
  implementation.
\end{enumerate}
\end{exercise}

\begin{exercise}[1.5]
Consider the problem of using a source of unbiased random bits to generate
samples from the set $S=\{0,\ldots,n-1\}$ such that the element $i$ is
chosen with probability $p_i$.  Show how to perform this sampling using
$O(\log n)$ random bits per sample, regardless of the values of $p_i$.  Use
the result from part~(c) of Problem~1.4.
\end{exercise}

\begin{exercise}[1.6]
Consider a sequence of $n$ flips of an unbiased coin.  Let $H_i$ denote the
absolute value of the excess of the number of HEADS over the number of TAILS
seen in the first $i$ flips.  Define $H=\max_i H_i$.  Show that
$\mathbb{E}[H_i]=\Theta(\sqrt{i})$, and that
$\mathbb{E}[H]=\Theta(\sqrt{n})$.
\end{exercise}

\begin{exercise}[1.7]
Suppose we choose a permutation $\pi$ of the ordered set
$N=\{1,2,\ldots,n\}$ uniformly at random from the space of all permutations
of $N$.  Let $L(\pi)$ denote the length of the longest increasing
subsequence in permutation~$\pi$.
\begin{enumerate}[label=(\alph*)]
\item For large $n$ and some positive constant $c$, prove that
  $\mathbb{E}[L(\pi)]\ge c\sqrt{n}$.
\item Is the bound in~(a) tight?
\end{enumerate}
\end{exercise}

\begin{exercise}[1.8]
Consider adapting the min-cut algorithm of Section~1.2.2 to the
problem of finding an $s$-$t$ min-cut in an undirected graph.  In this
problem, we are given an undirected graph $G$ together with two distinguished
vertices $s$ and $t$.  An $s$-$t$ cut is a set of edges whose removal from
$G$ disconnects $s$ from $t$; we seek an $s$-$t$ cut of minimum cardinality.
As the algorithm proceeds, the vertex $s$ may get amalgamated into a new
vertex as a result of an edge being contracted; we call this vertex the
$s$-vertex (initially the $s$-vertex is $s$ itself).  Similarly, we have a
$t$-vertex.  As we run the contraction algorithm, we ensure that we never
contract an edge between the $s$-vertex and the $t$-vertex.
\begin{enumerate}[label=(\alph*)]
\item Show that there are graphs on which the probability that this algorithm
  finds an $s$-$t$ min-cut is exponentially small.
\item How large can the number of $s$-$t$ min-cuts in an instance be?
\end{enumerate}
\end{exercise}

\begin{exercise}[1.9]
Consider the Find algorithm described in Section~\ref{sec:recurrence} for
selecting the $k$th smallest of a set $S$ of $n$ elements.  Show that the
algorithm finds the $k$th smallest element in $S$ in expected time $O(n)$.
\end{exercise}

\begin{exercise}[1.10]
Consider the setting of Example~1.1 (the Sailor Problem).  Show that the
probability that no sailor returns to her own cabin approaches $1/e$ as the
number of sailors grows large.
\end{exercise}

\begin{exercise}[1.11]
Verify the following inclusions:
$\mathrm{P}\subseteq\mathrm{RP}\subseteq\mathrm{NP}
  \subseteq\mathrm{PSPACE}\subseteq\mathrm{EXP}\subseteq\mathrm{NEXP}$.
It is not known whether these inclusions are strict.
\end{exercise}

\begin{exercise}[1.12]
Verify the following inclusions:
$\mathrm{RP}\subseteq\mathrm{BPP}\subseteq\mathrm{PP}$.
It is not known whether these inclusions are strict.
\end{exercise}

\begin{exercise}[1.13]
Show that $\mathrm{PP}=\mathrm{co\text{-}PP}$ and
$\mathrm{BPP}=\mathrm{co\text{-}BPP}$.
\end{exercise}

\begin{exercise}[1.14]
Show that $\mathrm{NP}\subseteq\mathrm{PP}\subseteq\mathrm{PSPACE}$.
\end{exercise}

\begin{exercise}[1.15]
(Due to K.-I.~Ko~[265].)  Show that $\mathrm{NP}\subseteq\mathrm{BPP}$
implies $\mathrm{NP}=\mathrm{RP}$.
\end{exercise}

\end{document}
